\clearpage
\VerdierPhysicalPageAnchor{0249}%
\begingroup\small\setlength{\parskip}{0pt}\setlength{\emergencystretch}{1.5em}
\setlength{\abovedisplayskip}{4pt plus 2pt minus 2pt}%
\setlength{\belowdisplayskip}{4pt plus 2pt minus 2pt}%

\VerdierSemanticAnchor{ch4d-4-6-9-first-differentials}
\noindent であり、第一微分は (4.1.9.1)
\[
\begin{aligned}
d_1^{p,q}\colon(\mathrm{C.E.})_1^{p,q}
  &\longrightarrow(\mathrm{C.E.})_1^{p+1,q},\\
d_1^{p,q}&=F^q(-d_2^{\bullet,p}).
\end{aligned}
\]
である。

\VerdierSemanticAnchor{ch4d-4-6-9-column-decomposition}
複体 $I^{\bullet,p}$ は Cartan--Eilenberg 入射複体であり (4.6.1)、したがって
(4.6.3) により $D(\mathcal A)$ に標準同型
\[
I^{\bullet,p}\xrightarrow{\sim}
\prod_{r\in\mathbb Z}H^r(I^{\bullet,p})[-r].
\]
がある。

\VerdierSemanticAnchor{ch4d-4-6-9-h1-components}
ここで (4.6.4.3) で導入した複体 $H_1^r(I^{\bullet\bullet})$ を用いる。
任意の整数 $r$ に対して
\[
H_1^r(I^{\bullet\bullet})^p=H^r(I^{\bullet,p}),
\]
であり、したがって $D(\mathcal A)$ に標準同型

\VerdierSemanticAnchor{ch4d-4-6-9-h1-decomposition}
\[
I^{\bullet,p}\xrightarrow{\sim}
\prod_{r\in\mathbb Z}H_1^r(I^{\bullet\bullet})^p[-r],
\]
がある。また、この同型のもとで複体の射
\[
d_2^{\bullet,p}\colon
I^{\bullet,p}\longrightarrow I^{\bullet,p+1}
\]
は射

\VerdierSemanticAnchor{ch4d-4-6-9-d2-decomposition}
\[
\prod_{r\in\mathbb Z}
d_{H_1^r(I^{\bullet\bullet})}^{p}[-r],
\]
と同型である。これにより初期項の最終的な表示

\VerdierSemanticAnchor{ch4d-4-6-9-1}
\[
(\mathrm{C.E.})_1^{p,q}
=F^q\left(
\prod_{r\in\mathbb Z}H_1^r(I^{\bullet\bullet})^p[-r]
\right),
\tag{4.6.9.1}
\]
および第一微分の最終的な表示

\VerdierSemanticAnchor{ch4d-4-6-9-2}
\[
d_1^{p,q}
=-F^q\left(
\prod_{r\in\mathbb Z}d_{H_1^r(I^{\bullet\bullet})}^{p}[-r]
\right).
\tag{4.6.9.2}
\]
を得る。

\VerdierSemanticAnchor{ch4d-4-6-9-3-prelude}
$F^{q,\bullet}(H_1^r(I^{\bullet\bullet}))$ を、複体
$H_1^r(I^{\bullet\bullet})$ の成分と微分に関手 $F^q$ を適用して得られる
$\mathcal B$ の対象の複体とする。公式 (4.6.9.1) と (4.6.9.2) は射

\VerdierSemanticAnchor{ch4d-4-6-9-3}
\[
(\mathrm{C.E.})_2^{p,q}\longrightarrow
\prod_{r\in\mathbb Z}
H_{\mathcal B}^{p}\left(
F^{q-r,\bullet}(H_1^r(I^{\bullet\bullet}))
\right).
\tag{4.6.9.3}
\]
を定める。
\endgroup
